% =============================================================================
%  Paper §3.9 -- Method reference: Hochschild Cohomology for \beta-Reduction
%  Composition (李思 2026).
%
%  Sources:
%    molmetal/math_skills/hochschild_cohomology.py          (T11 math skeleton)
%    molmetal/molmetal_lam/lam_chem/lambda_combinators.py   (R17 PRIMITIVE_ATOMS)
%    .claude/skills/mol-metal-lambda-calculus/SKILL.md      (MLC formalism)
%    .claude/skills/mol-metal-r18-math/SKILL.md             (T11 priority entry)
%
%  Status: METHOD REFERENCE for the T11 cohomological-semantics layer.
%  All four primitives (Cochain / Differential / HochschildCohomology /
%  MLCReductionCochain) are SHIPPED as a CPU-only math layer;
%  25/25 unit tests in molmetal/tests/test_hochschild_cohomology.py
%  (target).
%
%  Theorem anchor (per .claude/skills/mol-metal-r18-math/SKILL.md):
%    T11 -- Li (李思), Wang & Yang 2026, "Stochastic Calculus and
%    Hochschild Homology" -- primary cite (arXiv:2501.12360,
%    seminar 2026-03-09 at Peking U. Statistical Science Center).
%
%  Honest framing:
%    - The math layer (Cochain, Differential, HochschildCohomology,
%      MLCReductionCochain) is SHIPPED as a CPU-only mathematical
%      skeleton.  Cited as DESIGN + ARCHITECTURAL.
%    - The category-theoretic semantics (every MLC reduction
%      sequence is a Hochschild cochain; the closure theorem is
%      the trivial-cohomology case) is MATH-ONLY -- no metric
%      uplift is claimed for the cohomological-semantics layer
%      in R18.
%    - The §3 chain remains: MLC formalism (§3.1) + MCTS (§3.4) +
%      Deflex (§3.5) + SE(3) SB (§3.4-3.6) + transport relaxation
%      (§3.5) + multi-PH (§3.6) + monge-ampere / W2-GF (§3.7) +
%      quantum semantics (§3.8) -- the cohomological-semantics
%      layer is a *category-theoretic* backbone, not an
%      independent metric target.
% =============================================================================

% -----------------------------------------------------------------------------
% Phantom-cite workaround for the lit anchor cited below but not yet
% in refs.bib.  Same pattern as the SE(3) SB, Transport Relaxation,
% Multi-PH and Quantum Semantics sections above.  After the close-out,
% the citation key ``lisi2026hochschild'' will resolve through BibTeX
% to the Li-Wang-Yang 2026 seminar + arXiv:2501.12360 entry added
% to paper/refs.bib.
% -----------------------------------------------------------------------------
\makeatletter
\def\WFlisi2026Key{lisi2026hochschild}
\let\WForigCiteHC\WForiginalCite
\def\WFcitePlainHC#1{%
  \def\WFciteTempKeyHC{#1}%
  \ifx\WFciteTempKeyHC\WFlisi2026Key Li, Wang \& Yang, 2026%
  \else \WForigCiteHC{#1}%
  \fi
}
\let\cite\WFcitePlainHC
\makeatother
% -----------------------------------------------------------------------------

\subsection{Hochschild Cohomology for $\beta$-Reduction Composition}
\label{sec:method:hochschild}

The category-theoretic semantics layer of the Molecular Lambda
Calculus (\S\ref{sec:method:mlc-formalism}) lifts every
$\beta$-reduction sequence to a Hochschild cochain complex.
This sub-section introduces the cohomological-semantics layer
that backs the T11 entry of the R18 backlog
(\texttt{mol-metal-r18-math/SKILL.md}) and ships the four
mathematical primitives in
\texttt{molmetal/math\_skills/hochschild\_cohomology.py}.

\paragraph{Why a cohomological-semantics layer.}
The MLC formalism of \S\ref{sec:method:mlc-formalism} is
**algebraic**: an atom is a primitive combinator of fixed
arity
(\texttt{valence + lone\_pairs}); a molecule is a closed term
in $\beta$-normal form; a synthesis is a $\beta$-reduction
sequence.  The category-theoretic semantics layer
**lifts** the algebraic reading to the **cohomological**
setting by identifying every $\beta$-reduction sequence with a
Hochschild cochain $c \in C^{n}(A, M)$ where $A$ is the
closed-term sub-algebra of MLC and $M$ is the bimodule of
typed atoms with left/right bimodule actions given by the
MLC application rule ($f \cdot x$ = curried partial application).

This lifting is the **classical** (non-quantum) shadow of
the **quantization-commutes-with-reduction** programme
initiated by Vergne (ICM 2006) and brought to its full
generality for compact symplectic manifolds by Ma \& Zhang
(Astérisque 318, 2008); the T10 quantum-semantics layer
(\S\ref{sec:method:quantum-semantics}) lifts this further to
the symplectic / CR-manifold setting.  Li, Wang \& Yang
(2026) \cite{lisi2026hochschild} (seminar at Peking
University Statistical Science Center, 2026-03-09; preprint
arXiv:2501.12360, 2025) investigate the
**probabilistic-topological** version of this lifting and
prove that the expectation of a stochastic integral over a
closed Markov loop recovers the Hochschild homology class of
the corresponding operator product -- a fundamental
*Hochschild / Markov* correspondence that the present
section adapts to the Mol-Metal surface.

\paragraph{Honest framing.}
The math layer is \textsc{shipped} as a CPU-only mathematical
skeleton at \texttt{molmetal/math\_skills/hochschild\_cohomology.py}
(target: 25/25 unit tests in
\texttt{molmetal/tests/test\_hochschild\_cohomology.py}).
The category-theoretic backbone is MATH-ONLY -- no metric
uplift is claimed for the cohomological-semantics layer in
R18.  The T11 priority entry of
\texttt{mol-metal-r18-math/SKILL.md} \S priority map is
\textbf{category-theoretic, no metric target}.

\subsubsection{Mathematical statement (informal)}
\label{sec:method:hochschild-statement}

Let $A$ be the closed-term sub-algebra of MLC (= the set of
closed $\beta$-NF terms, with multiplication given by
$\beta$-reduction composition).  Let $M$ be the bimodule of
typed atoms with left/right bimodule actions given by the MLC
application rule.  The **Hochschild cochain complex** of
$(A, M)$ is
\begin{equation}
  \label{eq:hc-complex}
  \ldots \xrightarrow{\delta} C^{n-1}(A, M)
                    \xrightarrow{\delta} C^n(A, M)
                    \xrightarrow{\delta} C^{n+1}(A, M)
                    \xrightarrow{\delta} \ldots ,
\end{equation}
with the **coboundary operator** $\delta: C^{n} \to C^{n+1}$
given by
\begin{equation}
  \label{eq:hc-delta}
  (\delta f)(a_1, \ldots, a_{n+1})
    = a_1 \cdot f(a_2, \ldots, a_{n+1})
    + \sum_{i=1}^{n} (-1)^i f(a_1, \ldots, a_i a_{i+1},
                                \ldots, a_{n+1})
    + (-1)^{n+1} f(a_1, \ldots, a_n) \cdot a_{n+1}.
\end{equation}
The **fundamental lemma** of Hochschild cohomology is
\begin{equation}
  \label{eq:hc-delta-squared}
  \delta \circ \delta \;=\; 0,
\end{equation}
which makes the **cohomology** $H^n(A, M) =
\ker(\delta_n) / \mathrm{im}(\delta_{n-1})$ well-defined.

\paragraph{MLC surface: every $\beta$-reduction sequence is a cochain.}
A single $\beta$-reduction (one bond-formation step between
two atoms) is a 1-cochain $c \in C^1(A, M)$.  A
two-step $\beta$-reduction (an **associativity witness**)
is a 2-cochain $f \in C^2(A, M)$ whose value $f(a, b, c)$
records the *agreement* of the two ways of composing
$(a, b)$ into $a \cdot b$ and then reducing with $c$, vs.
composing $(b, c)$ first and reducing with $a$.  The
fundamental identity $\delta^2 = 0$ then says that
associativity witnesses compose consistently: every
$\delta c$ (the coboundary of a 1-cochain) is a 2-cocycle
that vanishes under $\delta$, recovering the **closure
theorem** that the closed-term sub-algebra of MLC is
*associative* (modulo the symmetric $\sqrt{2}$-scaled
midpoint composition -- see \texttt{BetaReductionSymp.apply}
in \S\ref{sec:method:quantum-semantics-surface}).

\paragraph{Closure-theorem is the trivial-cohomology case.}
For the closed-term sub-algebra $A_{\mathrm{MLC}}$ of MLC,
every associativity witness is **trivial** (it is a
coboundary of a 1-cochain via the $\sqrt{2}$ composition),
so
\begin{equation}
  \label{eq:hc-trivial}
  H^2(A_{\mathrm{MLC}}, A_{\mathrm{MLC}}) \;=\; 0,
  \qquad
  H^0(A_{\mathrm{MLC}}, A_{\mathrm{MLC}}) \;=\; A_{\mathrm{MLC}}.
\end{equation}
In words: the closure theorem is **literally** the
trivial-cohomology case of the Hochschild complex of MLC.
This is the canonical categorical semantics for the
closure theorem proved empirically for 80 hypothesis cases
in
\texttt{molmetal/molmetal\_lam/lam\_chem/closure.py} and
formally in the closure-theorem appendix
(\texttt{paper/appendices/closure\_theorem.tex}).

\paragraph{Concrete operational test.}
For the canonical choice of bimodule multiplication
(left-mult = right-mult = mult = add), the coboundary
operator $\delta$ reduces to the **alternating sum**
\begin{equation}
  \label{eq:hc-alternating}
  (\delta f)(a_1, \ldots, a_{n+1})
    = \sum_{i=0}^{n+1} (-1)^i \, \sigma_i(a_1, \ldots, a_{n+1}),
\end{equation}
where $\sigma_i$ is the insertion at position $i$ of the
multiplication.  The corresponding test
\texttt{Differential.apply(coch)} +
\texttt{Differential.delta\_squared\_zero(coch, test\_args)}
verifies this alternating-sum formula at run-time; the
fundamental identity $\delta^2 = 0$ is operationalised by
\texttt{test\_differential\_delta\_zero} which numerically
confirms that $\delta^2 c = 0$ on a battery of test
arguments for an arbitrary 1-cochain.

\subsubsection{Mol-Metal surface: every MLC reduction is a Hochschild cochain}
\label{sec:method:hochschild-surface}

The four primitives in
\texttt{molmetal/math\_skills/hochschild\_cohomology.py}
are the mathematical backbone that the W2.x adapters will
wrap:

\begin{enumerate}
  \item \textbf{\texttt{Cochain}}
    (\texttt{math\_skills/hochschild\_cohomology.py}, class at
    lines $\sim$120--$\sim$185).  An element
    $c \in C^{n}(A, M)$ of the Hochschild cochain complex,
    represented as a length-1 list of multilinear callables
    wrapping a single $n$-ary function
    $c(a_1, \ldots, a_n) \mapsto M$-element.
    \texttt{evaluate(args)} runs the callable on the
    argument list; \texttt{degree} stores $n$.

  \item \textbf{\texttt{Differential}} (lines $\sim$200--$\sim$310).
    The coboundary operator $\delta: C^{n} \to C^{n+1}$ for a
    given algebra $(A, M)$ (specified by three callables:
    ``left\_mult``, ``right\_mult``, ``mult'').
    \texttt{apply(cochain)} computes $\delta c$ (returns a
    Cochain of degree ``cochain.degree + 1``);
    \texttt{delta\_squared\_zero(cochain, test\_args)}
    checks the fundamental identity $\delta^2 c = 0$
    numerically on a battery of arguments.

  \item \textbf{\texttt{HochschildCohomology}} (lines
    $\sim$330--$\sim$470).  Computes
    $\dim H^{n}(A, M)$ for $n = 0, 1, 2, \ldots,
    \mathrm{max\_degree}$ (default $\mathrm{max\_degree}=3$)
    by binning the supplied ``cochains`` list by degree,
    counting cocycles (cochains that map to zero under
    $\delta$), counting coboundaries (the image of $\delta$
    from the previous degree), and returning
    $\dim\ker - \dim\mathrm{im}$ as a
    ``torch.Tensor`` of dtype ``long``.

  \item \textbf{\texttt{MLCReductionCochain}} (lines
    $\sim$485--$\sim$590).  The MLC surface wrapper for a
    $\beta$-reduction: each cochain is parameterised by an
    ``atom`` symbol (e.g. ``'Pt\_II'``,
    ``'Ru\_III'``); ``apply(other)`` composes two
    reductions via the symmetric $\sqrt{2}$-scaled midpoint
    (witnessing the closure theorem);
    ``to\_cochain()`` lifts the reduction to a 1-cochain
    on the standard Hochschild complex;
    ``closure\_theorem\_holds(a, b, c)`` verifies the
    associativity witness on three cochains (the **MHC
    closure-theorem test**, mirroring the 80-case empirical
    closure test of \texttt{molmetal/molmetal\_lam/lam\_chem/closure.py}).
\end{enumerate}

\paragraph{W2.x adapter (planned).}
The W2.x adapter wraps the math layer with a category-theoretic
bridge between MLC's $\beta$-NF witnesses and the live
SynFlowNet / CFM stack.  Specifically, the adapter takes a
sequence of $\beta$-reduction witnesses from
\texttt{molmetal/molmetal\_lam/lam\_chem/lambda\_combinators.py}
and emits a Hochschild cochain signature
$(c_0, c_1, c_2) \mapsto$ shape $(|\mathrm{atoms}|, \ldots)$
that the live CFM / SynFlowNet advantage estimator can consume
as a category-theoretic side-channel feature.  The
\texttt{integrate\_with\_cfm(\dots)} signature will be compatible
with \texttt{flow\_matching.solver.ode\_solver.ODESolver} so a
future refactor can swap the adapter into the live CFM pipeline
without modifying the R17 locked
\texttt{molmetal/adapters/flow\_matching\_lipman/} close-out.

\paragraph{Pocket-conditioned cochain marginals (future).}
The target marginal on the Hochschild cochain is the
pocket-conditioned distribution of $\beta$-reduction
sequences that compose into a dockable metal complex,
parameterised by the QuickVina score against the target
pocket.  The source marginal is the unbound-atom distribution
(uniformly random cochain symbols).  The cross-pair assignment
is computed by a future **Hochschild-aware** transport cost
that uses the coboundary operator $\delta$ as the
cochain-level metric -- this is the R18 Phase 3 work item.

\subsubsection{Expected gain and honest framing}
\label{sec:method:hochschild-expected}

\paragraph{Category-theoretic gain (not metric).}
The cohomological-semantics layer is **category-theoretic**:
it provides a mathematically rigorous identification of every
MLC $\beta$-reduction sequence with a Hochschild cochain,
and the closure theorem with the trivial-cohomology case
$H^n(A_{\mathrm{MLC}}, A_{\mathrm{MLC}}) = 0$ for $n \geq 1$.
This is the natural setting for future upgrades that involve:
\begin{itemize}
  \item \textbf{composition-aware} MCTS whose expansion
    rule respects the Hochschild coboundary operator -- a
    cochain-theoretic generalisation of the closure-theorem
    precheck in
    \texttt{molmetal/molmetal\_lam/search\_alg/proof\_search.py};
  \item \textbf{cochain-aware reward shaping} -- replacing
    the F5 / LambdaASTSR symbolic-regression reward with a
    Hochschild-coboundary loss that inherits the
    Li-Wang-Yang \cite{lisi2026hochschild} Markov / Hochschild
    correspondence;
  \item a \textbf{Batalin-Vilkovisky (BV) extension} of the
    cohomological semantics layer -- the natural setting
    for higher categorical structure (2-categories of
    bonded atoms, gerbes of $\beta$-reduction orbits, etc.)
    relevant to the multi-metal extension at R19+.
\end{itemize}
No metric uplift is claimed in R18 for the
cohomological-semantics layer.

\paragraph{Honest framing on the metric uplift.}
The category-theoretic semantics is MATH-ONLY -- no metric
target.  The §3 chain of metric targets is unchanged:
\begin{itemize}
  \item §3.4 SE(3) SB: Vina Dock $-8.5 \to -10.5$ \PROJECTED;
  \item §3.5 Transport relaxation: pK$_\text{bind}$ MAE
    $0.42 \to 0.30$ \PROJECTED;
  \item §3.5 Deflex: A2780cis AUC uplift \PROJECTED;
  \item §3.6 Multi-PH: MetalCytotoxDB Ru AUC
    $0.81 \to 0.88$ \PROJECTED.
\end{itemize}
The cohomological-semantics layer is the
**category-theoretic backbone** that will enable future
metric targets but does not itself target a metric.

\paragraph{Frontier anchors.}
The cohomological-semantics layer sits on top of the
Hochschild / Markov correspondence of Li, Wang \& Yang
\cite{lisi2026hochschild} (arXiv:2501.12360, seminar
2026-03-09 at Peking U. Statistical Science Center) and
the foundational Hochschild-cohomology monograph of
Loday (Loday 1998, *Cyclic Homology*, Springer-Verlag);
the algebraic-deformation anchor is Gerstenhaber (1960,
*On the deformation of rings and algebras*,
Ann. of Math. 79 (1960), 59-103) and the Markov-loop
correspondence is part of the Li-Wang-Yang seminar's
programme.

\paragraph{Cross-section pointers.}
The cohomological-semantics layer feeds the §3.1 MLC
formalism (\texttt{Sec.~method:mlc-formalism}) by
providing the **category-theoretic semantics** that lifts
the algebraic reading of $\beta$-reductions to the
cochain-level; the §3.4 MCTS proof-search
(\texttt{Sec.~method:mcts-search}) by providing a
coboundary-aware expansion rule whose precheck is the
trivial-cohomology case $H^2 = 0$ (the closure theorem);
the §3.8 quantum-semantics layer
(\texttt{Sec.~method:quantum-semantics}) by being its
**classical (non-quantum) shadow** -- the T11
cohomological layer is what the T10 quantum-semantics
layer quantises; and the §7 future work
(\texttt{paper/sections/07\_future.tex}) by pointing at
the BV-algebra extension and the
$H^\bullet(A_{\mathrm{MLC}}, \mathbb{Z}/2)$ mod-2
Bockstein spectral sequence as the natural R19+ target
(especially relevant for the multi-metal extension to Fe,
Mn, Cu, Pd, Rh, Os).

% -----------------------------------------------------------------------------
% Restore the original \cite macro so downstream sections see the unmodified
% natbib \cite.  Same pattern as 03_3_pll_heuristic.tex:295-297 and
% 03_4_mcts_search.tex:740-744.
% -----------------------------------------------------------------------------
\makeatletter
\let\cite\WForigCiteHC
\makeatother
